% Pista geo-24s-q6 · Geometria
% Procedència: PAU Catalunya 2024 · Convocatòria de setembre · Sèrie 3 · Qüestió 6
\documentclass[11pt,a4paper]{article}
\usepackage{pista}
\pistainfo{Catalunya 2024}{Convocatòria de setembre}{Sèrie 3}

\begin{document}

\section*{\textcolor{blauPista}{Qüestió 6}}

\noindent Considereu el punt $P = (1, 3, 0)$ i el pla $\pi$ d'equació $x + 2y - 2z = -7$.

\subsection*{\textit{a)} \normalfont Sigui $r$ la recta que és perpendicular a $\pi$ i passa per $P$. Calculeu el punt d'intersecció de $\pi$ amb $r$. \hspace*{\fill}\textnormal{[1~punt]}}

\pista{Com que la recta $r$ és perpendicular al pla $\pi$, el seu vector director és el vector normal de $\pi$. Les seves components són els coeficients de $x$, $y$ i $z$ en l'equació del pla.

\medskip
En primer lloc, escriu la recta $r$ en forma paramètrica, amb el punt $P$ i aquest vector director. A continuació, substitueix les expressions de $x$, $y$ i $z$ a l'equació de $\pi$: obtindràs una equació amb la variable $\lambda$. Un cop l'hagis resolta, substitueix el valor de $\lambda$ a l'equació de $r$ i obtindràs el punt d'intersecció que estaves buscant.}

\subsection*{\textit{b)} \normalfont Calculeu la distància $d$ del punt $P$ al pla $\pi$. \hspace*{\fill}\textnormal{[0,5~punts]}}

\pista{En primer lloc, escriu el pla en la forma \mbox{$Ax + By + Cz + D = 0$}, passant el $-7$ a l'altre costat:
\[
x + 2y - 2z + 7 = 0
\]
A continuació, aplica la fórmula de la distància d'un punt $(x_{0}, y_{0}, z_{0})$ a un pla:
\[
d(P, \pi) = \frac{|Ax_{0} + By_{0} + Cz_{0} + D|}{\sqrt{A^{2} + B^{2} + C^{2}}}
\]}

\subsection*{\textit{c)} \normalfont Calculeu l'equació d'un altre pla $\pi'$ que sigui paral·lel a $\pi$ i que també estigui a distància $d$ de $P$. \hspace*{\fill}\textnormal{[1~punt]}}

\pista{Fes un dibuix: si el pla $\pi'$ és paral·lel a $\pi$ i el punt $P$ és a la mateixa distància dels dos plans, $P$ queda just al mig, amb un pla a cada costat.

\medskip
Com que $\pi'$ és paral·lel a $\pi$, té el mateix vector normal. Per tant, la seva equació és d'aquesta forma:
\[
x + 2y - 2z + D' = 0
\]
A continuació, imposa que la distància de $P$ a $\pi'$ sigui $d$, amb la fórmula de l'apartat \textit{b)}.

\medskip
Obtindràs una equació amb un valor absolut. Recorda que una equació com $|X| = 5$ vol dir que $X = 5$ o bé $X = -5$. Una de les solucions és el pla $\pi$ de l'enunciat; l'altra solució correspon al pla $\pi'$ que busques.}

\end{document}
